\documentclass[11pt]{article}

\usepackage{amsmath, amssymb, amsthm}
\usepackage[margin=1in]{geometry}

\DeclareMathOperator{\Var}{Var}
\DeclareMathOperator{\Cov}{Cov}
\DeclareMathOperator{\Tr}{Tr}
\DeclareMathOperator{\disc}{disc}
\DeclareMathOperator*{\argmin}{arg\,min}
\DeclareMathOperator{\Bernoulli}{Bernoulli}

\title{A Mathematical Compendium of Algorithms, Machine Learning, Optimization, and Complexity}
\author{DeepSeek prompted by Riyad Derguini}
\date{}

\begin{document}

\maketitle

\textbf{Notation and Conventions.} Let \(\mathbb{R}\) denote the real numbers and \(\mathbb{R}^n\) the \(n\)-dimensional Euclidean space. Vectors are denoted by boldface lowercase letters (e.g., \(\mathbf{x}\)), and matrices by boldface uppercase letters (e.g., \(\mathbf{A}\)). The transpose of a vector or matrix is denoted by \(^\top\). The gradient of a scalar function \(f\) is \(\nabla f\), and its Hessian is \(\mathbf{H}\) or \(\nabla^2 f\). The indicator function of a set \(A\) is \(\mathbb{I}(A)\). For asymptotic notation, \(n\) denotes the input size. The natural logarithm is \(\log\), and \(\|\cdot\|_p\) denotes the \(L^p\) norm.

\hrulefill

\part{Complexity Theory and Algorithmic Analysis}

In algorithmic analysis, we abstract away hardware specifics to study resource scaling (time and space) in terms of the input size \(n\).

\section{Asymptotic Notation (Formal Definitions)}
Let \(f, g: \mathbb{N} \to \mathbb{R}_{>0}\) be functions.
\begin{itemize}
    \item \textbf{Big-O (\(O\)):} \(f(n) = O(g(n))\) if \(\exists c > 0, n_0 \in \mathbb{N}\) such that \(\forall n \ge n_0\), \(0 \le f(n) \le c \cdot g(n)\).
    \item \textbf{Big-Omega (\(\Omega\)):} \(f(n) = \Omega(g(n))\) if \(\exists c > 0, n_0 \in \mathbb{N}\) such that \(\forall n \ge n_0\), \(0 \le c \cdot g(n) \le f(n)\).
    \item \textbf{Big-Theta (\(\Theta\)):} \(f(n) = \Theta(g(n))\) if \(f(n) = O(g(n))\) and \(f(n) = \Omega(g(n))\).
    \item \textbf{Little-o (\(o\)):} \(f(n) = o(g(n))\) if \(\lim_{n \to \infty} \frac{f(n)}{g(n)} = 0\).
    \item \textbf{Little-omega (\(\omega\)):} \(f(n) = \omega(g(n))\) if \(\lim_{n \to \infty} \frac{f(n)}{g(n)} = \infty\).
\end{itemize}

\section{The Master Theorem}
For recurrences of the form \(T(n) = aT(n/b) + f(n)\), where \(a \ge 1, b > 1\), and \(f(n)\) is asymptotically positive:
\begin{enumerate}
    \item If \(f(n) = O(n^{\log_b a - \epsilon})\) for some \(\epsilon > 0\), then \(T(n) = \Theta(n^{\log_b a})\).
    \item If \(f(n) = \Theta(n^{\log_b a})\), then \(T(n) = \Theta(n^{\log_b a} \log n)\).
    \item If \(f(n) = \Omega(n^{\log_b a + \epsilon})\) for some \(\epsilon > 0\), and if \(a f(n/b) \le k f(n)\) for some constant \(k < 1\) and all sufficiently large \(n\), then \(T(n) = \Theta(f(n))\).
\end{enumerate}

\section{Complexity Classes}
\begin{itemize}
    \item \textbf{P:} Problems solvable by a deterministic Turing machine in \(O(n^k)\) time for some \(k\).
    \item \textbf{NP:} Problems where a given solution can be verified in polynomial time by a deterministic Turing machine.
    \item \textbf{NP-Complete:} A problem is NP-Complete if it is in NP and every problem in NP is polynomial-time reducible to it. (Cook-Levin Theorem establishes SAT as the first NP-Complete problem).
    \item \textbf{NP-Hard:} Problems at least as hard as NP-Complete problems, but not necessarily in NP.
    \item \textbf{co-NP:} Problems where a ``no'' certificate can be verified in polynomial time.
    \item \textbf{PSPACE:} Problems solvable in polynomial space.
    \item \textbf{BPP (Bounded-Error Probabilistic Polynomial Time):} Problems solvable by a randomized algorithm with error probability \(\le 1/3\) in polynomial time.
    \item \textbf{\#P (Sharp-P):} Counting problems associated with NP decision problems.
\end{itemize}

\section{Algorithmic Paradigms}
\begin{itemize}
    \item \textbf{Randomized Algorithms:}
    \begin{itemize}
        \item \emph{Las Vegas:} Always correct, but runtime is a random variable.
        \item \emph{Monte Carlo:} Bounded runtime, but correctness is probabilistic.
    \end{itemize}
    \item \textbf{Approximation Algorithms:} An algorithm is a \(\rho\)-approximation for an optimization problem if it runs in polynomial time and outputs a solution whose cost is within a factor \(\rho\) of the optimal cost.
    \item \textbf{Amortized Analysis:} Bounding the average time per operation over a worst-case sequence of operations.
\end{itemize}

\hrulefill

\part{Taylor Series and Mathematical Approximations}

Taylor series are fundamental to numerical optimization, algorithmic analysis, and function approximations in ML.

\section{The General Taylor Series}
The Taylor series of a function \(f(x)\) infinitely differentiable at a point \(a\) is:
\[
 f(x) = \sum_{n=0}^{\infty} \frac{f^{(n)}(a)}{n!} (x - a)^n 
\]
When \(a = 0\), this is called the Maclaurin series.

\section{The Multivariate Taylor Series}
For a multivariate function \(f: \mathbb{R}^n \to \mathbb{R}\), the second-order Taylor expansion around \(\mathbf{x}_0\) is:
\[
 f(\mathbf{x}) \approx f(\mathbf{x}_0) + \nabla f(\mathbf{x}_0)^\top (\mathbf{x} - \mathbf{x}_0) + \frac{1}{2} (\mathbf{x} - \mathbf{x}_0)^\top \mathbf{H}(\mathbf{x}_0) (\mathbf{x} - \mathbf{x}_0) 
\]
where \(\nabla f(\mathbf{x}_0)\) is the gradient vector and \(\mathbf{H}(\mathbf{x}_0)\) is the Hessian matrix. This is the basis for Newton's Method in optimization.

\section{Essential Maclaurin Series}
\begin{itemize}
    \item \( e^x = \sum_{n=0}^{\infty} \frac{x^n}{n!} \) for all \(x \in \mathbb{R}\).
    \item \( \log(1+x) = \sum_{n=1}^{\infty} (-1)^{n-1} \frac{x^n}{n} \) for \(x \in (-1, 1]\).
    \item \( \frac{1}{1-x} = \sum_{n=0}^{\infty} x^n \) for \(|x| < 1\).
    \item \( (1+x)^\alpha = \sum_{n=0}^{\infty} \binom{\alpha}{n} x^n \) for \(|x| < 1\).
    \item \( \sin(x) = \sum_{n=0}^{\infty} (-1)^n \frac{x^{2n+1}}{(2n+1)!} \).
    \item \( \cos(x) = \sum_{n=0}^{\infty} (-1)^n \frac{x^{2n}}{(2n)!} \).
    \item \textbf{Stirling's Approximation:} \( n! \sim \sqrt{2\pi n} \left(\frac{n}{e}\right)^n \), which implies \(\log(n!) = n \log n - n + O(\log n)\).
\end{itemize}

\hrulefill

\part{Probability and Statistics}

\section{Core Definitions and Rules}
\begin{itemize}
    \item \textbf{Bayes' Theorem:} \( P(A|B) = \frac{P(B|A)P(A)}{P(B)} \).
    \item \textbf{Law of Total Probability:} \( P(B) = \sum_i P(B|A_i)P(A_i) \).
    \item \textbf{Expectation:} \( \mathbb{E}[X] = \sum_x x P(X=x) \) or \( \int x f(x) dx \).
    \item \textbf{Variance:} \( \Var(X) = \mathbb{E}[X^2] - (\mathbb{E}[X])^2 \).
    \item \textbf{Covariance:} \( \Cov(X,Y) = \mathbb{E}[XY] - \mathbb{E}[X]\mathbb{E}[Y] \).
    \item \textbf{Tower Property:} \( \mathbb{E}[X] = \mathbb{E}[\mathbb{E}[X|Y]] \).
\end{itemize}

\section{Concentration Inequalities}
\begin{itemize}
    \item \textbf{Markov's Inequality:} For non-negative \(X\) and \(a > 0\), \( P(X \ge a) \le \frac{\mathbb{E}[X]}{a} \).
    \item \textbf{Chebyshev's Inequality:} \( P(|X - \mathbb{E}[X]| \ge a) \le \frac{\Var(X)}{a^2} \).
    \item \textbf{Chernoff Bounds:} For \(X = \sum X_i\) with independent Bernoulli variables, \( P(X \ge (1+\delta)\mu) \le \exp\left(-\frac{\delta^2 \mu}{2+\delta}\right) \).
    \item \textbf{Hoeffding's Inequality:} For independent random variables \(X_i \in [a_i, b_i]\), \( P(|\bar{X} - \mathbb{E}[\bar{X}]| \ge t) \le 2\exp\left(-\frac{2n^2t^2}{\sum (b_i-a_i)^2}\right) \).
    \item \textbf{Jensen's Inequality:} If \(g\) is convex, \( g(\mathbb{E}[X]) \le \mathbb{E}[g(X)] \).
\end{itemize}

\section{Key Distributions}
\begin{itemize}
    \item \textbf{Gaussian:} \( \mathcal{N}(x | \mu, \sigma^2) = \frac{1}{\sqrt{2\pi\sigma^2}} \exp\left(-\frac{(x-\mu)^2}{2\sigma^2}\right) \).
    \item \textbf{Multivariate Gaussian:} \( \mathcal{N}(\mathbf{x} | \boldsymbol{\mu}, \boldsymbol{\Sigma}) = \frac{1}{(2\pi)^{d/2}|\boldsymbol{\Sigma}|^{1/2}} \exp\left(-\frac{1}{2}(\mathbf{x}-\boldsymbol{\mu})^\top \boldsymbol{\Sigma}^{-1} (\mathbf{x}-\boldsymbol{\mu})\right) \).
    \item \textbf{Exponential Family:} \( p(x|\theta) = h(x) \exp(\eta(\theta)^\top T(x) - A(\theta)) \).
\end{itemize}

\section{Information Theory}
\begin{itemize}
    \item \textbf{Entropy:} \( H(X) = -\sum P(x) \log P(x) \).
    \item \textbf{Cross-Entropy:} \( H(P, Q) = -\sum P(x) \log Q(x) \).
    \item \textbf{KL Divergence:} \( D_{KL}(P || Q) = \sum P(x) \log \frac{P(x)}{Q(x)} \), with the identity \( H(P,Q) = H(P) + D_{KL}(P||Q) \).
\end{itemize}

\hrulefill

\part{Combinatorics and Discrepancy Theory}

\section{Fundamental Counting}
\begin{itemize}
    \item \textbf{Permutations:} \( P(n, k) = \frac{n!}{(n-k)!} \).
    \item \textbf{Combinations:} \( \binom{n}{k} = \frac{n!}{k!(n-k)!} \).
    \item \textbf{Vandermonde's Identity:} \( \binom{m+n}{r} = \sum_{k=0}^r \binom{m}{k}\binom{n}{r-k} \).
    \item \textbf{Stars and Bars:} \( \binom{n+k-1}{k-1} \) ways to distribute \(n\) indistinguishable balls into \(k\) distinguishable bins.
\end{itemize}

\section{The Probabilistic Method}
\begin{itemize}
    \item \textbf{Union Bound:} \( P(\cup_i A_i) \le \sum_i P(A_i) \).
    \item \textbf{Lovász Local Lemma (LLL):} Provides conditions under which the probability of avoiding a set of ``bad events'' with limited dependencies is strictly positive.
    \item \textbf{Method of Conditional Expectations:} A derandomization technique that fixes variables one by one to maintain the expected value.
\end{itemize}

\section{Discrepancy Theory}
\begin{itemize}
    \item \textbf{Standard Definition:} \( \disc(\mathcal{F}) = \min_{\chi \in \{\pm 1\}^n} \max_{S \in \mathcal{F}} |\sum_{x \in S} \chi(x)| \).
    \item \textbf{Matrix Form:} Minimizing \( \|A\mathbf{x}\|_\infty \), where \(A\) is the incidence matrix and \(\mathbf{x} \in \{\pm 1\}^n\).
    \item \textbf{Spencer's Theorem:} For \(n\) elements and \(m\) sets, \( \disc(\mathcal{F}) = O(\sqrt{m \log(n/m)}) \).
\end{itemize}

\hrulefill

\part{Optimization and Gradient Methods}

\section{Convexity}
\begin{itemize}
    \item \textbf{First-Order Condition:} \( f(y) \ge f(x) + \nabla f(x)^\top (y - x) \).
    \item \textbf{Second-Order Condition:} The Hessian is positive semi-definite, \( \mathbf{H} \succeq 0 \).
\end{itemize}

\section{Gradient Descent (GD)}
\begin{itemize}
    \item \textbf{Update Rule:} \( \mathbf{w}_{t+1} = \mathbf{w}_t - \eta \nabla \mathcal{L}(\mathbf{w}_t) \).
    \item \textbf{Convergence (Smooth Convex):} \( O(1/\epsilon) \) iterations to reach \(\epsilon\)-accuracy.
    \item \textbf{Stochastic Gradient Descent (SGD):} \( \mathbf{w}_{t+1} = \mathbf{w}_t - \eta \nabla \mathcal{L}_i(\mathbf{w}_t) \), where \(\nabla \mathcal{L}_i\) is an unbiased estimate of the true gradient.
\end{itemize}

\section{Newton's Method}
\begin{itemize}
    \item \textbf{Update Rule:} \( \mathbf{w}_{t+1} = \mathbf{w}_t - [\mathbf{H}\mathcal{L}(\mathbf{w}_t)]^{-1} \nabla \mathcal{L}(\mathbf{w}_t) \).
    \item \textbf{Convergence:} Quadratic convergence near the minimum.
\end{itemize}

\section{Regularization}
\begin{itemize}
    \item \textbf{L2 (Ridge):} \( \Omega(\mathbf{w}) = \frac{\lambda}{2} \|\mathbf{w}\|_2^2 \).
    \item \textbf{L1 (Lasso):} \( \Omega(\mathbf{w}) = \lambda \|\mathbf{w}\|_1 \), which promotes sparsity.
    \item \textbf{Lagrangian Formulation:} For \(\min f(x)\) subject to \( g_i(x) \le 0 \), the Lagrangian is \( \mathcal{L}(x, \alpha) = f(x) + \sum \alpha_i g_i(x) \).
    \item \textbf{KKT Conditions:} Stationarity, Primal Feasibility, Dual Feasibility, and Complementary Slackness.
\end{itemize}

\hrulefill

\part{Machine Learning Specifics}

\section{Bias-Variance Decomposition}
\[
 \mathbb{E}[(y - \hat{f}(x))^2] = (\text{Bias}[\hat{f}(x)])^2 + \Var[\hat{f}(x)] + \sigma^2_{\text{irreducible}} 
\]
\begin{itemize}
    \item \textbf{Overfitting:} Characterized by high variance and low bias.
    \item \textbf{Underfitting:} Characterized by high bias and low variance.
\end{itemize}

\section{Gradients of Common Loss Functions}
\begin{itemize}
    \item \textbf{MSE:} \( \mathcal{L} = \frac{1}{2}(y - \hat{y})^2 \); gradient with respect to \(\hat{y}\) is \( \hat{y} - y \).
    \item \textbf{Binary Cross-Entropy:} \( \mathcal{L} = -[y \log \hat{y} + (1-y)\log(1-\hat{y})] \); gradient with respect to logits is \( \hat{y} - y \).
    \item \textbf{Softmax + Categorical Cross-Entropy:} Gradient with respect to logits is \( \hat{\mathbf{y}} - \mathbf{y} \).
\end{itemize}

\section{Matrix Calculus Identities}
\begin{itemize}
    \item \( \frac{\partial}{\partial \mathbf{x}} (\mathbf{a}^\top \mathbf{x}) = \mathbf{a} \).
    \item \( \frac{\partial}{\partial \mathbf{x}} (\mathbf{x}^\top \mathbf{A} \mathbf{x}) = (\mathbf{A} + \mathbf{A}^\top)\mathbf{x} \).
    \item \( \frac{\partial}{\partial \mathbf{X}} \Tr(\mathbf{A}\mathbf{X}) = \mathbf{A}^\top \).
    \item \( \frac{\partial}{\partial \mathbf{X}} \log |\det(\mathbf{X})| = (\mathbf{X}^{-1})^\top \).
\end{itemize}

\section{Singular Value Decomposition (SVD)}
For any matrix \(\mathbf{A}\), the SVD is
\[
 \mathbf{A} = \mathbf{U} \boldsymbol{\Sigma} \mathbf{V}^\top. 
\]
This is the foundation of Principal Component Analysis (PCA) and low-rank approximation (Eckart-Young-Mirsky theorem).

\end{document}
